\section{One hereditary FULL-C arm suffices against an arbitrary other arm}
This is an unbounded enlargement of the ORIGINAL coefficient theorem.
It does not assert general C or solve the case where neither arm has
a hereditary C certificate. There are ZERO new Lean reports: the
weighted-sum, monotone-excess, closed-window and source/payment kernels
are the already verified inputs. The class induction and composition
assembly below are written proofs, not an end-to-end formalization.

\paragraph{Marked classes and exact statements.}
Let $\mathcal C$ be ANY final-part-deletion-closed class of normalized
actual increasing arms, with FULL C for every nonrectangular member.
Include the empty arm and all rectangular prefixes that occur on deletion.
The normalization is $X=1^\alpha U$, $\alpha\ge1$, with $U$ sorted and
all its parts at least two. For $A=F_X$ put $G_k=\sum_{i=0}^k A_i$,
$D_n=(A_n-A_{n-1})^+$, $D_0=D_1=0$. FULL C is
\[
 G_{k+1}D_n\ge G_kD_{n+1}
 \qquad(0\le k\le n-3,\ D_nD_{n+1}>0).                 \tag{C}
\]
Write $\mathcal U$ for ALL actual increasing arms. A marked pair is
an unordered pair in
\[
          \mathcal P=(\mathcal C\times\mathcal U)
                         \cup(\mathcal U\times\mathcal C).
\]
The unmarked member has no capacity, count, initial-run, first-missing,
nonbinary-weight, or C restriction. Marking persists under deletion
of EITHER arm. The proof does not replace the hereditary hypothesis
by C of the single current marked arm.

\begin{theorem}
For every marked pair $(X,Y)$, with $q_T=z^{\operatorname{wt}T}F_T(1/z)$,
the product $p=q_Xq_Y$ has FULL A2:
\[
                  \bigl(2^n(p_n-p_{n-1})^+\bigr)_{n\ge0}
                    \text{ is unimodal.}                         \tag{1}
\]
For every strict peak $M>\max(X,Y)$, EACH of the ORIGINAL polynomials
\[
\begin{split}
 F_X+F_Y+(S_M-1)F_XF_Y,\qquad
 F_X+(S_M-1)F_XF_Y,\qquad F_Y+(S_M-1)F_XF_Y
\end{split}                                                       \tag{2}
\]
is coefficient unimodal. No amplification is performed.
Unlimited legal appends on the unmarked arm are already allowed.
In addition, ONE arbitrary legal append $X=Zt$ on a marked arm
$Z\in\mathcal C$ is permitted, provided
$\operatorname{wt}(Zt)\ge\operatorname{wt}Y$. Its new C is not required.
\end{theorem}

\subsection{The two uses of C each need only the marked arm}
First suppose both arms are nonrectangular. At every rank where both
forward arm polynomials still rise, orient the marked arm as $A$ and
write $B$ for the other arm. With $E_B=1-z+(2z-1)B$, the actual source
increments $w_j=\Delta E_B(j)$ are nonnegative throughout this prefix,
and $w_j=0$ for $j<3$. Consequently
\[
 (E_BA)_n=\sum_{j=3}^n w_jG_{n-j}.
\]
Summing (C) with these nonnegative weights gives the source/slope
ratio comparison. This is exactly the proof in
\texttt{clipped\_forward\_ratio\_frontier.tex}; it uses C only for $A$.
After either forward factor stops rising, the symmetric actual source
identity makes $2(F_XF_Y)_{n-1}-(F_XF_Y)_n$ nonnegative forever.
Thus coefficient tilt A1 holds for every nonrectangular marked pair.

The same one-owner weighted comparison, together with monotonicity of
the source convolution, gives the nonnegative-current nonrecovery
in \texttt{C\_to\_coupled\_S\_rising\_window.tex}. Its statement already
explicitly requires C of EITHER arm, rather than both arms. This supplies
the two-forward-rising part of the proper CLOSED source comparison.
Neither step differentiates the unmarked arm's slope or invokes its C.

\subsection{The other phases and the closed-window bridge}
For the two-forward-nonrising phase use the unconditional actual
first-half and reverse half-tail resources, as in
\texttt{complete\_mixed\_source\_unrestricted\_C\_original\_criterion\_proof.tex}.
For the mixed phase use the COMPLETE proper CLOSED mixed-source
comparison of that proof. Its small-first-missing source/head payments,
the complete 10--12 certificate and the unbounded 13--15 and 16-plus
payments concern arbitrary actual owners and opposites. The owner
need not be the marked arm. FULL OLD A2 supplies its head budget;
no C of the mixed owner or opposite is used in those payments.
The one-non-one-part opposite and the negative-rank unit birth are
included in that complete theorem.

Together these inputs prove EVERY proper CLOSED source comparison,
including equality $E_n=E_k$. The complete bridge in
\texttt{coupled\_source\_window\_A2\_and\_kappa\_nine\_original\_proof.tex}
therefore gives FULL A2 for EVERY uniform window $[t]p$, assuming
FULL OLD A2 and A1. The bridge explicitly retains its first head
boundary, zero positive-difference support and equality case. No
whole raw-source-window unimodality is substituted for this bridge.

\subsection{Count induction for all marked pairs}
Induct on the total number of non-one parts and orient the heavier
arm as $X$. If either arm is rectangular, the unconditional rectangular
base in \texttt{binary\_ternary\_A2\_original\_family\_proof.tex} applies
with an arbitrary actual other arm. If the heavier arm has exactly
one non-one part, the heavier-one-part theorem in
\texttt{deficit\_cumulative\_ratios.tex} also allows an arbitrary
nonrectangular other arm. Neither base requires its C.

Otherwise write the heavier arm as $X=Zt$. The pair $(Z,Y)$ is again
marked: deleting an unmarked final part leaves the marked member
untouched; deleting a marked final part leaves its prefix in
$\mathcal C$. Its smaller total count supplies FULL OLD A2. If both
old arms are nonrectangular, the preceding two subsections give FULL
A2 for every $[t](q_Zq_Y)$. Rectangular old arms are the same bases.
The exact new product is
\[
 q_Xq_Y=[t](q_Zq_Y)+z^a q_Y,\qquad a=\operatorname{wt}X.
\]
The strengthened actual ordinary midpoint theorem gives its first
mode $\sigma<a$, including equal nonrectangular weights. The terminal
summand starts at $a$, above $\sigma$, so positive differences through
$\sigma$ are exactly those of the window. Their weighted prefix is
unimodal. All later actual positive differences vanish. This proves
(1), with FULL old A2, not a truncated inductive substitute.

For the extra append on a marked arm, its old pair $(Z,Y)$ has the
just-proved FULL A2, A1 and every closed comparison. The same argument
works for the arbitrary legal $t$ under $a\ge\operatorname{wt}Y$.
No additional step of that argument asks for C of $Zt$.

\subsection{Original core and both relatives, with all closure checks}
Adjoin to $\mathcal P$ the exceptional pairs $(Zt,Y)$ specified in the
theorem, and permit arm exchange. This enlarged family is closed
under final deletion. In an exceptional pair, deleting its exceptional
jump returns to $\mathcal P$; deleting from $Y$ preserves the heavier
weight inequality. In an ordinary marked pair either deletion preserves
the mark or its marked prefix. Removing leading ones stays within the
rectangular bases when the non-one word is empty.

This family is also closed under the lower-peak stripping used in
\texttt{short\_prefix\_reflection\_join.tex} and
\texttt{inductive\_signed\_window.tex}. For an ordinary marked pair,
stripping all occurrences of a current largest arm capacity is a
sequence of final deletions. In an exceptional pair, if that largest
capacity is $t$, strip the exceptional jump first and return to a
marked pair. If the largest capacity exceeds $t$, its occurrences
are all in $Y$; stripping them preserves the exceptional heavier arm.
Repeat until the lower strict-peak or rectangular base is reached.

Thus the least-peak/least-relative argument applies literally: FULL
A2 supplies its weighted positive-difference input, proper deletion
supplies its already-unimodal relative input, and all reflection and
lower-peak joins stay inside the enlarged family. Both relatives in
(2) follow, and heavier-first core-relative equivalence gives the core.
This proves the theorem at EVERY strict peak, not just one chosen peak.

\paragraph{Explicit newly covered unbounded original families.}
The marked arm may have at most FIVE arbitrary non-one parts; or any
number of parts at most FOUR; or at most TWO arbitrary nonbinary parts
after an arbitrary binary prefix; or any fixed nonbinary word of total
weight at most 32. These are hereditary FULL-C seeds from the frozen
certificates. One may further append ANY FINITE NUMBER of successive
dominant capacities $T_i\ge2(W_0+\sum_{v<i}T_v)$, with the established
empty-seed and first-jump conditions. The other arm is COMPLETELY
ARBITRARY, including long 6/7 arms that violate every fixed-count-offset
ordinary log-concavity claim. These examples use the appropriate
hereditary closure of each seed, rather than only its current C.

\subsection{A stronger negative-current source budget from the same mark}
Let $P=AB$, $e_N=2P_{N-1}-P_N$, $w_A=\Delta E_A$, $w_B=\Delta E_B$.
If $A$ has FULL C and $e_N<0$, then
\[
 (E_AE_B)_N\le w_A(N),\qquad
 e_N-2e_{N-1}\ge w_B(N)\ge0.                            \tag{3}
\]
Here $A$ need not have hereditary C: only this local conclusion uses
current C. To prove it, negative $e_N$ forces both factors still to
rise, and the one-owner A1 argument above gives $e_{N-1}\le0$.
Put $d=D_{N-1}>0$, $S=(E_BA)_{N-1}$, so $S-d=e_{N-1}\le0$.
For $3\le j\le N-1$, set $k=N-1-j$. The corresponding C inequality is
\[
 w_A(N)G_k=(2d-D_N)G_k
                 \ge d(2G_k-G_{k+1})
                 =d\sum_{i=0}^{k+1}E_A(i).
\]
Sum against the nonnegative $w_B(j)$; the new boundary terms have
zero cumulative source. This proves
\[
 d(E_AE_B)_N\le w_A(N)S\le w_A(N)d.
\]
Cancellation proves the first bound in (3). Finally the exact identity
\[
 w_A+w_B-E_AE_B=(1-z)^2+(1-2z)e
\]
at rank $N\ge3$ proves the second. The initial negative rank $N=2$
has $e_2=-1$, $e_1=0$, so it MUST be excluded from (3); its source
identity has the nonzero $(1-z)^2$ correction. For $N=3$ the sum
is empty, and the same displayed argument is valid. If BOTH arms have
C, the first bound holds with $\min(w_A(N),w_B(N))$.

\paragraph{Remaining general problem.}
The two-unknown-arm situation is still OPEN. The theorem proves an
original coefficient statement whenever ONE arm has hereditary C;
it does not prove that every actual arm has such a certificate.
Finite actual-model checks in the accompanying audit are QA only.
